% ============================================================
% 哲学论文学术报告 Handout LaTeX 模板 (Markdown极简版)
% 编译方式: xelatex (需编译两次以获得正确交叉引用)
% ============================================================
\documentclass[10pt, a4paper]{article}

% -------------------- 基础宏包 --------------------
\usepackage[UTF8]{ctex}                    % 中文支持
\usepackage[margin=2cm]{geometry}          % 页边距
\usepackage{amsmath, amssymb, amsthm}      % 数学与逻辑符号
\usepackage{booktabs}                       % 专业表格线
\usepackage{tabularx}                       % 弹性宽度表格
\usepackage{enumitem}                       % 自定义列表
\usepackage{hyperref}                       % 超链接
\usepackage{tikz}                           % 绘图
\usetikzlibrary{shapes.geometric, arrows.meta, positioning}
\usepackage{titlesec}                       % 标题格式
\usepackage{fancyhdr}                       % 页眉页脚
\usepackage{multicol}                       % 多栏排版
\usepackage{setspace}                       % 行距

% -------------------- 极简排版设置 --------------------
\setstretch{1.15}                          % 行距
\setlength{\parskip}{0.4em}                % 段落间距
\setlength{\parindent}{0pt}                % 首行缩进
\setlength{\abovedisplayskip}{4pt}
\setlength{\belowdisplayskip}{4pt}

% -------------------- 列表设置 --------------------
\setlist[itemize]{leftmargin=1.5em, itemsep=2pt, parsep=0pt, topsep=3pt}
\setlist[enumerate]{leftmargin=2em, itemsep=2pt, parsep=0pt, topsep=3pt}

% -------------------- 页眉页脚 --------------------
\pagestyle{fancy}
\fancyhf{}
\fancyhead[L]{\footnotesize\textit{\thetitle}}
\fancyhead[R]{\footnotesize\theauthor}
\fancyfoot[C]{\small\thepage}
\renewcommand{\headrulewidth}{0.3pt}

% -------------------- 标题格式 (极简) --------------------
\titleformat{\section}{\large\bfseries\sffamily}{\thesection}{0.6em}{}
\titleformat{\subsection}{\normalsize\bfseries}{\thesubsection}{0.5em}{}
\titleformat{\subsubsection}{\normalsize\itshape}{\thesubsubsection}{0.5em}{}
\titlespacing*{\section}{0pt}{1.2em}{0.4em}
\titlespacing*{\subsection}{0pt}{0.8em}{0.3em}
\titlespacing*{\subsubsection}{0pt}{0.6em}{0.2em}

% -------------------- 超链接 (黑白) --------------------
\hypersetup{
  colorlinks=true,
  linkcolor=black,
  citecolor=black,
  urlcolor=black,
  bookmarks=true,
  pdfstartview=FitH
}

% -------------------- 自定义命令 --------------------

% 英文术语标注
\newcommand{\engterm}[1]{\textup{\textsf{\small(#1)}}}

% 原文页码标注
\newcommand{\pagemark}[1]{\textsuperscript{\scriptsize[#1]}}

% 高亮文本
\newcommand{\highlight}[1]{\textbf{#1}}
\newcommand{\coreclaim}[1]{\underline{\textbf{#1}}}

% 脚注式评论标记
\newcommand{\remark}[1]{\textit{\small[#1]}}

% 双语分隔符
\newcommand{\zh}[1]{\\{\small\textmd{#1}}}

% 分隔线
\newcommand{\secsep}{\par\vspace{0.3em}\hrule\vspace{0.3em}\par}

% 定义环境（纯文本风格）
\newenvironment{deflist}{
  \begin{description}[style=nextline, leftmargin=1em, itemsep=0.3em]
}{
  \end{description}
}

% -------------------- 文档信息 --------------------
\title{}
\author{}
\date{}

% ============================================================
%  以下为 Handout 正文模板 — 请替换为实际内容
% ============================================================

\begin{document}

% ==================== ENGLISH VERSION ====================

{\centering
  {\Large\bfseries\sffamily Paper Title}\\[3pt]
  {\normalsize Author Name (Year) · \textit{Journal/Source}}\\[4pt]
  Presenter: Name · Course: Course Name · Date: YYYY.MM.DD\par
}

\vspace{0.3em}
\hrule
\vspace{0.5em}

\section*{Thesis}

\textbf{Core claim:} One-sentence summary of the main thesis.\pagemark{page}

Brief elaboration of the core claim in 2-3 sentences.

\section*{Background}

Context and theoretical background. Key assumptions challenged:
\begin{itemize}
  \item Assumption 1: Description\pagemark{page}
  \item Assumption 2: Description\pagemark{page}
\end{itemize}

\section*{Main Arguments}

\subsection*{Argument 1: Name}

\textbf{Setup:} Description of the thought experiment or scenario.\pagemark{page}

\textbf{Formalization:}
\begin{align*}
  &\text{P1: Premise 1} \\
  &\text{P2: Premise 2} \\
  &\therefore \text{C: Conclusion}
\end{align*}

\textbf{Evaluation:} Brief assessment of argument strength.

\section*{Key Concepts}

\begin{center}
\begin{tabular}{@{}lll@{}}
\toprule
\textbf{Term} & \textbf{Definition} & \textbf{Ref} \\
\midrule
Term 1 & Definition & \pagemark{page} \\
Term 2 & Definition & \pagemark{page} \\
\bottomrule
\end{tabular}
\end{center}

\section*{Commentary}

Assessment of the argument:\\[0.3em]
\begin{tabular}{@{}p{3.5cm}p{10cm}@{}}
\textbf{Strengths:} & List main strengths \\[0.2em]
\textbf{Weaknesses:} & List potential objections \\[0.2em]
\textbf{Significance:} & Theoretical impact
\end{tabular}

\newpage

% ==================== CHINESE VERSION ====================

{\centering
  {\Large\bfseries\sffamily 论文标题}\\[3pt]
  {\normalsize 作者名 (年份) · \textit{期刊/来源}}\\[4pt]
  汇报人：姓名 · 课程：课程名 · 日期：YYYY.MM.DD\par
}

\vspace{0.3em}
\hrule
\vspace{0.5em}

\section*{核心论题}

\textbf{核心主张：}用一句话概括中心论点。\pagemark{页码}

用2-3句话简要阐述核心主张。

\section*{理论背景}

背景和理论语境。被挑战的关键假设：
\begin{itemize}
  \item 假设1：描述\pagemark{页码}
  \item 假设2：描述\pagemark{页码}
\end{itemize}

\section*{核心论证}

\subsection*{论证1：名称}

\textbf{设置：}思想实验或场景的描述。\pagemark{页码}

\textbf{形式化：}
\begin{align*}
  &\text{P1: 前提1} \\
  &\text{P2: 前提2} \\
  &\therefore \text{C: 结论}
\end{align*}

\textbf{评价：}对论证强度的简要评估。

\section*{关键概念}

\begin{center}
\begin{tabular}{@{}lll@{}}
\toprule
\textbf{术语} & \textbf{定义} & \textbf{页码} \\
\midrule
术语1 & 定义 & \pagemark{页码} \\
术语2 & 定义 & \pagemark{页码} \\
\bottomrule
\end{tabular}
\end{center}

\section*{评论}

对论证的评价：\\[0.3em]
\begin{tabular}{@{}p{3.5cm}p{10cm}@{}}
\textbf{优点：} & 列出主要优点 \\[0.2em]
\textbf{弱点：} & 列出可能的异议 \\[0.2em]
\textbf{意义：} & 理论影响
\end{tabular}

\end{document}
